\subsection{Spectral Theory of Structured Operators}

% **Role:** structure revealed by invariant subspaces and eigenbehavior.
%
% Normal, self-adjoint, positive, unitary operators; Schur theorem; spectral theorem; spectral resolutions.

Structured spectral theory studies operators whose geometry and spectrum align. Normal, self-adjoint, positive, and unitary maps admit orthogonal decompositions and spectral resolutions.

\begin{definition} \label{def:normal-operator} \hlt{Normal Operator}\\
An operator $T\in\LL(\mathcal{V})$ on a finite-dim inner-product space is \hlt{normal} if it commutes with its adjoint: $TT^*=T^*T$.\\
A square matrix $A$ is \hlt{normal} if $A^*A=AA^*$.\\
Every self-adjoint operator, every Hermitian matrix, and every unitary matrix is normal.
\end{definition}

\begin{theorem} \label{thm:normal-iff-equal-adjoint-norms} \hlt{Normal Operators and Equal Adjoint Norms}\\
An operator $T\in\LL(\mathcal{V})$ is normal if and only if $\norm{T\mathbf{v}}=\norm{T^*\mathbf{v}}$ for every $\mathbf{v}\in\mathcal{V}$.
\begin{proof}
The operator $A=T^*T-TT^*$ is self-adjoint.\\
For every $\mathbf{v}\in\mathcal{V}$, $\innerprod{A\mathbf{v},\mathbf{v}} =\innerprod{T^*T\mathbf{v},\mathbf{v}}-\innerprod{TT^*\mathbf{v},\mathbf{v}} =\norm{T\mathbf{v}}^2-\norm{T^*\mathbf{v}}^2$.\\
Thus $T$ is normal exactly when $A=0$, which by Proposition \ref{prop:self-adjoint-zero-quadratic-form-implies-zero} is equivalent to the displayed quadratic form being zero for every $\mathbf{v}$. This is equivalent to $\norm{T\mathbf{v}}=\norm{T^*\mathbf{v}}$ for every $\mathbf{v}$.\\
\hlt{Structural Commentary}: Normality is measured by the self-adjoint commutator $T^*T-TT^*$.\\
\hlt{Key Insight}: A normal operator and its adjoint move every vector the same distance.
\end{proof}
\end{theorem}

\begin{theorem} \label{thm:normal-operator-adjoint-eigenvectors} \hlt{Normal Operators Share Eigenvectors with Their Adjoints}\\
Let $T\in\LL(\mathcal{V})$ be normal. If $T\mathbf{v}=\lambda\mathbf{v}$ for some $\mathbf{v}\neq\mathbf{0}$, then $T^*\mathbf{v}=\overline{\lambda}\mathbf{v}$.
\begin{proof}
The operator $T-\lambda I_{\mathcal{V}}$ is normal by Proposition \ref{prop:basic-adjoint-identities} and scalar multiples of $I_{\mathcal{V}}$ commute with $T$ and $T^*$. Applying Theorem \ref{thm:normal-iff-equal-adjoint-norms} to $T-\lambda I_{\mathcal{V}}$ gives $0=\norm{(T-\lambda I_{\mathcal{V}})\mathbf{v}} =\norm{(T-\lambda I_{\mathcal{V}})^*\mathbf{v}} =\norm{(T^*-\overline{\lambda}I_{\mathcal{V}})\mathbf{v}}$.\\
\hlt{Structural Commentary}: For normal operators, kernels of $T-\lambda I$ and $T^*-\overline{\lambda}I$ coincide.\\
\hlt{Key Insight}: Normality makes eigenvector information visible from both $T$ and $T^*$.
\end{proof}
\end{theorem}

\begin{figure}[H]
  \centering
  \begin{tikzpicture}[
  >=Stealth,
  every node/.style={font=\small},
  axis/.style={gray!55, line width=0.6pt},
  vec/.style={->, blue!75!black, line width=0.95pt},
  adj/.style={->, teal!70!black, line width=0.9pt},
  note/.style={draw=gray!45, rounded corners=3pt, align=center, inner sep=4pt}
]
  \coordinate (O) at (0,0);
  \coordinate (v) at (2.15,0.8);
  \coordinate (Tv) at (3.0,1.12);
  \coordinate (Tsv) at (2.55,0.95);

  \draw[axis] (-0.45,0) -- (3.55,0);
  \draw[axis] (0,-0.35) -- (0,1.85);
  \draw[vec] (O) -- (v) node[midway, above left] {$\mathbf{v}$};
  \draw[vec] (O) -- (Tv) node[above right] {$T\mathbf{v}=\lambda\mathbf{v}$};
  \draw[adj] (O) -- (Tsv) node[below right] {$T^*\mathbf{v}=\overline{\lambda}\mathbf{v}$};
  \draw[gray!60, dashed] (v) -- (Tv);
  \draw[gray!60, dashed] (v) -- (Tsv);

  \node[note] at (1.55,-0.95)
    {normality aligns the eigendirection\\for both $T$ and $T^*$};
\end{tikzpicture}

  \caption{For normal $T$, $T\mathbf{v}=\lambda\mathbf{v}$ implies $T^*\mathbf{v}=\overline{\lambda}\mathbf{v}$.}
  \label{fig:normal-operator-eigenvectors}
\end{figure}

\begin{theorem} \label{thm:normal-eigenvectors-distinct-eigenvalues-orthogonal} \hlt{Orthogonal Eigenvectors for Normal Operators}\\
Let $T\in\LL(\mathcal{V})$ be normal. Eigenvectors of $T$ corresponding to distinct eigenvalues are orthogonal.
\begin{proof}
Suppose $T\mathbf{u}=\alpha\mathbf{u}$ and $T\mathbf{v}=\beta\mathbf{v}$ with $\alpha\neq\beta$. By Theorem \ref{thm:normal-operator-adjoint-eigenvectors}, $T^*\mathbf{v}=\overline{\beta}\mathbf{v}$.\\
Hence $\alpha\innerprod{\mathbf{u},\mathbf{v}} =\innerprod{T\mathbf{u},\mathbf{v}} =\innerprod{\mathbf{u},T^*\mathbf{v}} =\innerprod{\mathbf{u},\overline{\beta}\mathbf{v}} =\beta\innerprod{\mathbf{u},\mathbf{v}}$. Since $\alpha-\beta\neq0$, $\innerprod{\mathbf{u},\mathbf{v}}=0$.\\
\hlt{Structural Commentary}: Normality lets the operator move from one eigenvector to the other through the adjoint.\\
\hlt{Key Insight}: Distinct eigenvalue scalings are incompatible unless the inner product is zero.
\end{proof}
\end{theorem}

\begin{theorem} \label{thm:self-adjoint-eigenvalues-real} \hlt{Eigenvalues of Self-Adjoint Operators Are Real}\\
Every eigenvalue of a self-adjoint operator on a finite-dimensional inner-product space is real.
\begin{proof}
Let $T=T^*$ and $T\mathbf{v}=\lambda\mathbf{v}$ with $\mathbf{v}\neq\mathbf{0}$. Then $\lambda\norm{\mathbf{v}}^2 =\innerprod{\lambda\mathbf{v},\mathbf{v}} =\innerprod{T\mathbf{v},\mathbf{v}} =\innerprod{\mathbf{v},T\mathbf{v}} =\innerprod{\mathbf{v},\lambda\mathbf{v}} =\overline{\lambda}\norm{\mathbf{v}}^2$.\\
Since $\norm{\mathbf{v}}^2>0$, $\lambda=\overline{\lambda}$, so $\lambda\in\mathbb{R}$.\\
\hlt{Structural Commentary}: Self-adjointness moves $T$ between inner-product slots without changing the value.\\
\hlt{Key Insight}: An eigenvalue must equal its conjugate because the eigenvector has positive squared norm.
\end{proof}
\end{theorem}

\begin{proposition} \label{prop:invertible-quadratic-self-adjoint} \hlt{Invertible Quadratic Expressions}\\
Let $T\in\LL(\mathcal{V})$ be self-adjoint on a finite-dimensional inner-product space.\\
If $b,c\in\mathbb{R}$ and $b^2<4c$, then $T^2+bT+cI_{\mathcal{V}}$ is invertible.
\begin{proof}
Let $S=T^2+bT+cI_{\mathcal{V}}$. For $\mathbf{v}\neq\mathbf{0}$,
\begin{align}
\innerprod{S\mathbf{v},\mathbf{v}} &=\norm{T\mathbf{v}}^2+b\innerprod{T\mathbf{v},\mathbf{v}}+c\norm{\mathbf{v}}^2 \geq \norm{T\mathbf{v}}^2-\abs{b}\norm{T\mathbf{v}}\norm{\mathbf{v}}+c\norm{\mathbf{v}}^2 \nonumber \\
&=\left(\norm{T\mathbf{v}}-\frac{\abs{b}}{2}\norm{\mathbf{v}}\right)^2+\left(c-\frac{b^2}{4}\right)\norm{\mathbf{v}}^2>0 \nonumber
\end{align}
Thus $S\mathbf{v}\neq\mathbf{0}$ for every nonzero $\mathbf{v}$, so $\ker(S)=\{\mathbf{0}\}$. In finite dimensions, injectivity implies invertibility.\\
\hlt{Structural Commentary}: Negative discriminant keeps this self-adjoint quadratic positive.\\
\hlt{Key Insight}: Self-adjointness converts $\innerprod{T^2\mathbf{v},\mathbf{v}}$ into $\norm{T\mathbf{v}}^2$, so completing the square controls the kernel.
\end{proof}
\end{proposition}

\begin{lemma} \label{lem:self-adjoint-operators-have-eigenvalues} \hlt{Self-Adjoint Operators Have Eigenvalues}\\
Let $\mathcal{V}\neq\{\mathbf{0}\}$ be a finite-dimensional real inner-product space. Every self-adjoint $T\in\LL(\mathcal{V})$ has an eigenvalue.
\begin{proof}
Choose $\mathbf{v}\neq\mathbf{0}$ and set $n=\dim(\mathcal{V})$.\\
The list $\mathbf{v},T\mathbf{v},\ldots,T^n\mathbf{v}$ is linearly dependent, so some nonzero real polynomial $p$ satisfies $p(T)\mathbf{v}=\mathbf{0}$.\\
Factor $p$ over $\mathbb{R}$ as $p(x)=a\prod\limits_{j=1}^{M}(x^2+b_jx+c_j)\prod\limits_{k=1}^{m}(x-\lambda_k)$, where $a\neq0$, $b_j^2<4c_j$, $\lambda_k\in\mathbb{R}$.\\
By Proposition \ref{prop:invertible-quadratic-self-adjoint}, each quadratic factor
$Q_j(T)\defeq T^2+b_jT+c_jI_{\mathcal{V}}$ is invertible. Let
$R(T)=\prod_j Q_j(T)$ and $L(T)=\prod_k(T-\lambda_kI_{\mathcal{V}})$. The factors commute because they are polynomials in $T$, and $R(T)$ is invertible.\\
If $m=0$, then $p(T)=aR(T)$ would be invertible, contradicting $p(T)\mathbf{v}=\mathbf{0}$ with $\mathbf{v}\neq\mathbf{0}$. Hence $m\geq1$.\\
Since $\mathbf{0}=p(T)\mathbf{v}=aR(T)L(T)\mathbf{v}$ and $aR(T)$ is invertible, $L(T)\mathbf{v}=\mathbf{0}$. Thus $L(T)$ has nontrivial kernel.\\
If every $T-\lambda_kI_{\mathcal{V}}$ were injective, their product $L(T)$ would be injective, a contradiction. Therefore some $T-\lambda_kI_{\mathcal{V}}$ is not injective, and $\lambda_k$ is an eigenvalue of $T$.\\
\hlt{Structural Commentary}: An annihilating polynomial remains; quadratic factors are invertible.\\
\hlt{Key Insight}: The nontrivial kernel must come from a real linear factor, hence from an eigenvalue.
\end{proof}
\end{lemma}

\begin{proposition} \label{prop:self-adjoint-invariant-subspaces} \hlt{Self-Adjoint Operators and Invariant Subspaces}\\
Let $T\in\LL(\mathcal{V})$ be self-adjoint, and let $\mathcal{U}$ be a subspace of $\mathcal{V}$ invariant under $T$. Then:
\begin{enumerate}[label=\roman*.]
\setlength{\itemsep}{0pt}
\item $\mathcal{U}^{\perp}$ is invariant under $T$;
\item $T|_{\mathcal{U}}\in\LL(\mathcal{U})$ is self-adjoint;
\item $T|_{\mathcal{U}^{\perp}}\in\LL(\mathcal{U}^{\perp})$ is self-adjoint.
\end{enumerate}
\begin{proof}
\begin{enumerate}[label=(\roman*)]
\setlength{\itemsep}{0pt}
\item
If $\mathbf{v}\in\mathcal{U}^{\perp}$ and $\mathbf{u}\in\mathcal{U}$, then $\innerprod{T\mathbf{v},\mathbf{u}}=\innerprod{\mathbf{v},T\mathbf{u}}=0$ because $T\mathbf{u}\in\mathcal{U}$. Hence $T\mathbf{v}\in\mathcal{U}^{\perp}$.
\item For $\mathbf{u},\mathbf{v}\in\mathcal{U}$, $\innerprod{(T|_{\mathcal{U}})\mathbf{u},\mathbf{v}} =\innerprod{T\mathbf{u},\mathbf{v}} =\innerprod{\mathbf{u},T\mathbf{v}} =\innerprod{\mathbf{u},(T|_{\mathcal{U}})\mathbf{v}}$. Thus $T|_{\mathcal{U}}$ is self-adjoint.
\item By (i), $\mathcal{U}^{\perp}$ is invariant, so (ii)'s argument applies to $T|_{\mathcal{U}^{\perp}}$.
\end{enumerate}
\hlt{Structural Commentary}: Self-adjointness preserves orthogonal complements of invariant subspaces.\\
\hlt{Key Insight}: Move $T$ across the inner product; invariance of $\mathcal{U}$ then forces orthogonality to persist.
\end{proof}
\end{proposition}

\begin{theorem} \label{thm:schur-triangularization} \hlt{Schur Triangularisation}\\
Let $\mathcal{V}$ be a finite-dimensional complex inner-product space and let $T\in\LL(\mathcal{V})$. Then $T$ has an upper-triangular matrix with respect to some orthonormal basis of $\mathcal{V}$.
\begin{proof}
Proceed by induction on $n=\dim(\mathcal{V})$. The cases $n=0$ and $n=1$ are immediate.\\
Suppose $n\geq2$ and the result holds for smaller dimensions. By Theorem \ref{thm:complex-operator-has-eigenvalue}, choose a unit eigenvector $\mathbf{e}_1$ of $T$, say $T\mathbf{e}_1=\lambda\mathbf{e}_1$. Let $\mathcal{U}=\operatorname{span}(\mathbf{e}_1)$ and define $S\in\LL(\mathcal{U}^{\perp})$ by $S\mathbf{w}=P_{\mathcal{U}^{\perp}}(T\mathbf{w})$.\\
By induction, choose an orthonormal basis $(\mathbf{e}_2,\ldots,\mathbf{e}_n)$ of $\mathcal{U}^{\perp}$ with respect to which $S$ has an upper-triangular matrix. Then $(\mathbf{e}_1,\ldots,\mathbf{e}_n)$ is an orthonormal basis of $\mathcal{V}$.\\
For $j\geq2$, the coordinates of $T\mathbf{e}_j$ along $\mathbf{e}_2,\ldots,\mathbf{e}_n$ are exactly the coordinates of $S\mathbf{e}_j$, hence have no entries below row $j$. Also $T\mathbf{e}_1=\lambda\mathbf{e}_1$. Therefore $[T]_{\gamma\leftarrow\gamma}$ is upper triangular for $\gamma=(\mathbf{e}_1,\ldots,\mathbf{e}_n)$.\\
\hlt{Structural Commentary}: Schur triangularisation builds an eigenline first, then triangularises the projected action on its orthogonal complement.\\
\hlt{Key Insight}: The lower-right block is controlled by the compression $P_{\mathcal{U}^{\perp}}T|_{\mathcal{U}^{\perp}}$.
\end{proof}
\end{theorem}

\begin{theorem} \label{thm:matrix-schur-triangularization-ordered} \hlt{Matrix Schur Triangularisation with Ordered Diagonal}\\
Let $A\in\operatorname{Mat}(n,n,\mathbb{C})$ have eigenvalues $\lambda_1,\ldots,\lambda_n$ in a prescribed order. If $A\mathbf{x}=\lambda_1\mathbf{x}$ and $\norm{\mathbf{x}}_2=1$, then there is a unitary $U=[\mathbf{x}\ U_2]$ such that $U^*AU$ is upper triangular with diagonal $\lambda_1,\ldots,\lambda_n$.\\
If $A$ is real, all $\lambda_i\in\mathbb{R}$, and $\mathbf{x}\in\mathbb{R}^n$, then $U$ may be chosen real orthogonal.
\begin{proof}
Induct on $n$. The case $n=1$ is immediate. Extend $\mathbf{x}$ to an orthonormal basis and set $V=[\mathbf{x}\ V_2]$.\\
Since $A\mathbf{x}=\lambda_1\mathbf{x}$, the first column of $V^*AV$ is $\lambda_1\mathbf{e}_1$, so $V^*AV=
\begin{bmatrix}
\lambda_1 & *\\
0 & A'
\end{bmatrix}$.\\
Block triangularity gives $p_{V^*AV}(z)=(z-\lambda_1)p_{A'}(z)$, so $A'$ has eigenvalues $\lambda_2,\ldots,\lambda_n$, counted with algebraic multiplicity. By induction, choose unitary $W$ with $W^*A'W=T'$ upper triangular and diagonal $\lambda_2,\ldots,\lambda_n$.\\
Then $U\defeq V(I_1\oplus W)$ has first column $\mathbf{x}$ and $U^*AU=
\begin{bmatrix}
\lambda_1 & *\\
0 & T'
\end{bmatrix}$.\\
The real case uses real orthonormal extensions; the block $A'$ is real, and the same characteristic-polynomial factorization leaves the real eigenvalue list $\lambda_2,\ldots,\lambda_n$ for the induction step.\\
\hlt{Structural Commentary}: Isolate the chosen eigenline, then repeat on the lower-right block.\\
\hlt{Key Insight}: The diagonal order is controlled by which eigenline is peeled off first.
\end{proof}
\end{theorem}

\begin{figure}[H]
  \centering
  \begin{tikzpicture}[
  >=Stealth,
  every node/.style={font=\small},
  block/.style={draw=gray!55, fill=gray!5, rounded corners=2pt, inner sep=6pt, align=center, minimum height=1.0cm},
  focus/.style={draw=blue!65!black, fill=blue!7, rounded corners=2pt, inner sep=6pt, align=center, minimum height=1.0cm},
  arr/.style={->, line width=0.75pt, blue!70!black}
]
  \node[focus] (line) at (0,0)
    {$A\mathbf{x}=\lambda_1\mathbf{x}$};
  \node[block] (block) at (3.3,0)
    {$V^*AV=\begin{bmatrix}\lambda_1&*\\0&A'\end{bmatrix}$};
  \node[block] (tri) at (6.9,0)
    {$U^*AU=T$};

  \draw[arr] (line) -- (block);
  \draw[arr] (block) -- (tri);
\end{tikzpicture}

  \caption{Schur induction: isolate an eigenline, then triangularise the compressed block.}
  \label{fig:schur-ordered-triangularization}
\end{figure}

\begin{corollary} \label{cor:eigenvalues-of-inverse-matrix} \hlt{Eigenvalues of the Inverse}\\
Let $A\in\operatorname{Mat}(n,n,\mathbb{C})$ be invertible with eigenvalues $\lambda_1,\ldots,\lambda_n$, counted with algebraic multiplicity. Then $A^{-1}$ has eigenvalues $\lambda_1^{-1},\ldots,\lambda_n^{-1}$, counted with algebraic multiplicity.
\begin{proof}
Take a Schur form $A=UTU^*$ with diagonal $\lambda_1,\ldots,\lambda_n$.\\
Then $A^{-1}=UT^{-1}U^*$, and $T^{-1}$ is upper triangular with diagonal $\lambda_1^{-1},\ldots,\lambda_n^{-1}$.\\
\hlt{Structural Commentary}: Invert the Schur form and read the diagonal.\\
\hlt{Key Insight}: The inverse spectrum is the reciprocal spectrum.
\end{proof}
\end{corollary}

\begin{theorem} \label{thm:complex-spectral-theorem} \hlt{Complex Spectral Theorem}\\
Let $\mathcal{V}$ be a finite-dimensional complex inner-product space and let $T\in\LL(\mathcal{V})$. The following are equivalent:
\begin{enumerate}[label=\roman*.]
\setlength{\itemsep}{0pt}
\item $T$ is normal;
\item $\mathcal{V}$ has an orthonormal basis consisting of eigenvectors of $T$;
\item $T$ has a diagonal matrix with respect to some orthonormal basis of $\mathcal{V}$.
\end{enumerate}
\begin{proof}
\begin{enumerate}[label=(\roman*).]
\setlength{\itemsep}{0pt}
\item[$(i)$]$\Rightarrow(ii)$
By Schur triangularisation (Theorem \ref{thm:schur-triangularization}), choose an orthonormal basis $\gamma=(\mathbf{e}_1,\ldots,\mathbf{e}_n)$ in which $A=[T]_{\gamma\leftarrow\gamma}=[a_{ij}]$ is upper triangular. Show by induction on $k$ that the $k$th row has no entries to the right of the diagonal. Assume this has been proved for rows $1,\ldots,k-1$.\\
Upper triangularity gives $T\mathbf{e}_k=\sum\limits_{i\leq k}a_{ik}\mathbf{e}_i$ and $T^*\mathbf{e}_k=\sum\limits_{j\geq k}\overline{a_{kj}}\mathbf{e}_j$.\\
The induction hypothesis gives $a_{ik}=0$ for $i<k$, so $\norm{T\mathbf{e}_k}^2=\abs{a_{kk}}^2$.\\
Since $T$ is normal, Theorem \ref{thm:normal-iff-equal-adjoint-norms} gives $\abs{a_{kk}}^2 =\norm{T\mathbf{e}_k}^2 =\norm{T^*\mathbf{e}_k}^2 =\abs{a_{kk}}^2+\sum\limits_{j>k}\abs{a_{kj}}^2$.\\
Thus each $a_{kj}$ with $j>k$ is zero. Induction makes $A$ diagonal, so $\gamma$ is an orthonormal basis of eigenvectors.
\item[$(ii)$]$\Rightarrow(iii)$
In an orthonormal eigenbasis, the matrix of $T$ is diagonal.
\item[$(iii)$]$\Rightarrow(i)$
If $[T]$ is diagonal in an orthonormal basis, then $[T^*]=[T]^*$ is diagonal as well.\\
Diagonal matrices commute, so $TT^*=T^*T$; hence $T$ is normal.
\end{enumerate}
\hlt{Structural Commentary}: Schur gives triangular form; normality kills the upper entries.\\
\hlt{Key Insight}: Equal norms of $T\mathbf{e}_j$ and $T^*\mathbf{e}_j$ force each row tail to vanish.
\end{proof}
\end{theorem}

\begin{figure}[H]
  \centering
  \begin{tikzpicture}[
  >=Stealth,
  every node/.style={font=\small},
  block/.style={draw=gray!55, rounded corners=3pt, inner sep=6pt, align=center},
  arrow/.style={->, line width=0.8pt, blue!70!black},
  faint/.style={gray!35, line width=0.7pt}
]
  \node[block] (schur) at (0,0)
    {$[T]_{\gamma}$ upper triangular\\$\begin{bmatrix}
      \lambda_1 & * & *\\
      0 & \lambda_2 & *\\
      0 & 0 & \lambda_3
    \end{bmatrix}$};
  \node[block] (normal) at (4.4,0)
    {normality\\$\norm{T\mathbf{e}_j}=\norm{T^*\mathbf{e}_j}$};
  \node[block] (diag) at (8.9,0)
    {$[T]_{\gamma}$ diagonal\\$\begin{bmatrix}
      \lambda_1 & 0 & 0\\
      0 & \lambda_2 & 0\\
      0 & 0 & \lambda_3
    \end{bmatrix}$};

  \draw[arrow] (schur) -- (normal);
  \draw[arrow] (normal) -- (diag);

  \draw[faint] (-1.1,0.58) -- (-0.53,0.58);
  \draw[faint] (-0.53,0.58) -- (-0.53,-0.18);
\end{tikzpicture}

\caption{Normality turns Schur triangular form diagonal by forcing row tails to vanish.}
  \label{fig:spectral-theorem-triangular-to-diagonal}
\end{figure}

\begin{theorem} \label{thm:real-spectral-theorem} \hlt{Real Spectral Theorem}\\
Let $\mathcal{V}$ be a finite-dimensional real inner-product space and let $T\in\LL(\mathcal{V})$. The following are equivalent:
\begin{enumerate}[label=\roman*.]
\setlength{\itemsep}{0pt}
\item $T$ is self-adjoint;
\item $\mathcal{V}$ has an orthonormal basis consisting of eigenvectors of $T$;
\item $T$ has a diagonal matrix with respect to some orthonormal basis of $\mathcal{V}$.
\end{enumerate}
\begin{proof}
\begin{enumerate}[label=(\roman*).]
\setlength{\itemsep}{0pt}
\item[$(i)$]$\Rightarrow (ii)$
Assume $T=T^*$ and induct on $n=\dim(\mathcal{V})$. For $n=0$ there is nothing to prove, and for $n=1$ every nonzero vector is an eigenvector after normalization.\\
Let $n\geq2$. Lemma \ref{lem:self-adjoint-operators-have-eigenvalues} gives a unit eigenvector $\mathbf{u}$, and set $\mathcal{U}=\operatorname{span}(\mathbf{u})$. The subspace $\mathcal{U}$ is $T$-invariant. By Proposition \ref{prop:self-adjoint-invariant-subspaces}, $\mathcal{U}^{\perp}$ is $T$-invariant and the restriction $T|_{\mathcal{U}^{\perp}}$ is self-adjoint.\\
Since $\dim(\mathcal{U}^{\perp})=n-1$, the induction hypothesis gives an orthonormal eigenbasis of $\mathcal{U}^{\perp}$.\\
Adjoining $\mathbf{u}$ gives an orthonormal eigenbasis of $\mathcal{V}$ by Theorem \ref{thm:orthogonal-direct-sum-decomposition}.
\item[$(ii)$]$\Rightarrow (iii)$
In an orthonormal eigenbasis, the matrix of $T$ is diagonal.
\item[$(iii)$]$\Rightarrow (i)$
In an orthonormal basis where $[T]$ is real diagonal, $[T^*]=[T]^t=[T]$. Hence $T=T^*$.
\end{enumerate}
\hlt{Structural Commentary}: Split off one eigendirection, then restrict to its orthogonal complement.\\
\hlt{Key Insight}: One eigenvector reduces the dimension.
\end{proof}
\end{theorem}

\begin{perspective} \label{per:normal-operators-coordinate-free-diagonal} \hlt{Normal Means Diagonal in Orthogonal Coordinates}\\
Complex spectral theorem states that normal operators are precisely the operators whose full structure is visible in an orthonormal eigenbasis. This is why normality is the natural coordinate-free version of being diagonal.
\end{perspective}

\begin{proposition} \label{prop:normal-polynomial-kernel-range-properties} \hlt{Basic Spectral Properties of Normal Operators}\\
Let $\mathcal{V}$ be a finite-dimensional complex inner-product space, let $T\in\LL(\mathcal{V})$ be normal, and let $p\in\mathbb{C}[z]$.
\begin{enumerate}[label=(\roman*)]
\setlength{\itemsep}{0pt}
\item $p(T)$ is normal and commutes with $T$.
\item $T^*$ is normal and $E_{\lambda}(T)=E_{\overline{\lambda}}(T^*)$ for every $\lambda\in\mathbb{C}$.
\item $\ker T=\ker T^*$ and $\operatorname{range}(T)=\operatorname{range}(T^*)$.
\end{enumerate}
\begin{proof}
Choose an orthonormal eigenbasis $\mathbf{e}_1,\ldots,\mathbf{e}_n$ with $T\mathbf{e}_j=\mu_j\mathbf{e}_j$.
\begin{enumerate}[label=(\roman*)]
\setlength{\itemsep}{0pt}
\item $p(T)\mathbf{e}_j=p(\mu_j)\mathbf{e}_j$, $p(T)$ is diagonal in this basis, hence normal. It commutes with $T$ by Proposition \ref{prop:operator-polynomial-products}.
\item Since $T^*\mathbf{e}_j=\overline{\mu_j}\mathbf{e}_j$, $T^*$ is diagonal in the same basis, hence normal. Also $T\mathbf{v}=\lambda\mathbf{v}$ iff $T^*\mathbf{v}=\overline{\lambda}\mathbf{v}$.
\item Part (ii) with $\lambda=0$ gives $\ker T=\ker T^*$. Both ranges are $\operatorname{span}\{\mathbf{e}_j:\mu_j\neq0\}$.
\end{enumerate}
\hlt{Structural Commentary}: Normality gives one orthonormal basis that diagonalises both $T$ and $T^*$.\\
\hlt{Key Insight}: Kernel and range are read from the nonzero diagonal entries in that basis.
\end{proof}
\end{proposition}

\begin{theorem} \label{thm:spectral-resolution-normal-operator} \hlt{Spectral Resolution of a Normal Operator}\\
Let $\mathcal{V}$ be a finite-dimensional complex inner-product space and let $\mu_1,\ldots,\mu_d$ be distinct scalars.\\
For $T\in\LL(\mathcal{V})$, consider:
\begin{enumerate}[label=(\roman*)]
\setlength{\itemsep}{0pt}
\item $T$ is normal and $\operatorname{spec}(T)\subseteq\{\mu_1,\ldots,\mu_d\}$.
\item $\exists$ pairwise orthogonal projections $P_1,\ldots,P_d\in\LL(\mathcal{V})$ such that $P_iP_j=0$ for $i\neq j$,\\
$P_1+\cdots+P_d=I_{\mathcal{V}}$, and $T=\mu_1P_1+\cdots+\mu_dP_d$
\item Any projections satisfying (ii) obey $\operatorname{range}(P_j)=E_{\mu_j}(T)$ whenever $\mu_j\in\operatorname{spec}(T)$.
\end{enumerate}
Then (i) and (ii) are equivalent, and any projections satisfying (ii) satisfy (iii).
\begin{proof}
\begin{enumerate}[label=(\roman*).]
\setlength{\itemsep}{0pt}
\item[$(i)$] $\Rightarrow (ii)$. By Theorem \ref{thm:complex-spectral-theorem}, $\mathcal{V}=\bigoplus\limits_{\lambda\in\operatorname{spec}(T)}E_{\lambda}(T)$ as an orthogonal direct sum.\\
Let $P_j$ be the orthogonal projection onto $E_{\mu_j}(T)$, and set $P_j=0$ when $\mu_j\notin\operatorname{spec}(T)$. Since $\operatorname{spec}(T)\subseteq\{\mu_1,\ldots,\mu_d\}$, these projections are pairwise orthogonal, sum to $I_{\mathcal{V}}$, and satisfy $T=\sum_j\mu_jP_j$.
\item[$(ii)$] $\Rightarrow (i)$. Since $P_j^*=P_j$ and $P_iP_j=0$ for $i\neq j$, $T^*=\sum\limits_{j=1}^{d}\overline{\mu_j}P_j$ and $T^*T=\sum\limits_{j=1}^{d}\abs{\mu_j}^2P_j=TT^*$.\\
Thus $T$ is normal. If $\lambda\in\operatorname{spec}(T)$ and $\mathbf{0}\neq\mathbf{x}\in E_{\lambda}(T)$, then $0=(T-\lambda I_{\mathcal{V}})\mathbf{x} =\sum\limits_{j=1}^{d}(\mu_j-\lambda)P_j\mathbf{x}$.\\
The summands lie in the pairwise orthogonal ranges $\operatorname{range}(P_j)$, hence $0=\sum\limits_{j=1}^{d}\abs{\mu_j-\lambda}^2\norm{P_j\mathbf{x}}^2$.\\
Because $\sum_jP_j=I_{\mathcal{V}}$, at least one $P_k\mathbf{x}$ is nonzero, so $\mu_k=\lambda$. Thus $\operatorname{spec}(T)\subseteq\{\mu_1,\ldots,\mu_d\}$.
\item[$(ii)$] $\Rightarrow (iii)$. For $\mathbf{x}\in\mathcal{V}$, $\mathbf{x} =\sum\limits_{i=1}^{d}P_i\mathbf{x}$ and $(T-\mu_jI_{\mathcal{V}})\mathbf{x} =\sum\limits_{i\neq j}(\mu_i-\mu_j)P_i\mathbf{x}$.\\
The summands are orthogonal, so $\norm{(T-\mu_jI_{\mathcal{V}})\mathbf{x}}^2 =\sum\limits_{i\neq j}\abs{\mu_i-\mu_j}^2\norm{P_i\mathbf{x}}^2$,
which vanishes exactly when $P_i\mathbf{x}=0$ for all $i\neq j$, equivalently $\mathbf{x}=P_j\mathbf{x}$. Thus $E_{\mu_j}(T)=\operatorname{range}(P_j)$.
\end{enumerate}
\hlt{Structural Commentary}: A normal operator decomposes into scalar actions on orthogonal eigenspaces.\\
\hlt{Key Insight}: Spectral projections isolate eigenspaces; the operator is the weighted sum of those projections.
\end{proof}
\end{theorem}

\begin{figure}[H]
  \centering
  \begin{tikzpicture}[
  >=Stealth,
  every node/.style={font=\small},
  block/.style={draw=gray!55, rounded corners=3pt, inner xsep=8pt, inner ysep=6pt, align=center},
  arrow/.style={->, line width=0.8pt, blue!70!black}
]
  \node[block] (space) at (0,0) {$\mathcal{V}$};
  \node[block] (decomp) at (3.2,0) {$E_{\mu_1}(T)\oplus\cdots\oplus E_{\mu_d}(T)$};
  \node[block] (operator) at (6.35,0) {$T$};

  \draw[arrow] (space) -- (decomp);
  \draw[arrow] (decomp) -- (operator);
\end{tikzpicture}

  \caption{Spectral resolution: $\mathcal{V}=\bigoplus_j\operatorname{range}(P_j)$, $P_iP_j=0$ for $i\neq j$, $\sum_jP_j=I_{\mathcal{V}}$, and $T=\sum_j\mu_jP_j$.}
  \label{fig:spectral-resolution-projectors}
\end{figure}

\begin{corollary} \label{cor:normal-spectral-resolution-functional-calculus} \hlt{Functional Calculus from Spectral Resolution}\\
Let $T=\sum_{j=1}^{d}\mu_jP_j$ be the spectral resolution of a normal operator and let $f\in\mathbb{C}[z]$. Then $f(T)=\sum\limits_{j=1}^{d}f(\mu_j)P_j$.\\
In particular, for each $j$ there is a polynomial $p_j$ such that $p_j(T)=P_j$.
\begin{proof}
Since $P_iP_j=0$ for $i\neq j$ and $P_j^2=P_j$, induction gives $T^m=\sum_j\mu_j^mP_j$ for every $m\geq0$.\\
Linearity then gives the displayed formula for every polynomial $f$.\\
For the final assertion, choose the Lagrange polynomial $p_j$ satisfying $p_j(\mu_i)=\delta_{ij}$. Then $p_j(T)=P_j$.\\
\hlt{Structural Commentary}: Orthogonal spectral projections multiply like coordinate idempotents.\\
\hlt{Key Insight}: Polynomial functions of a normal operator act by scalar functions on each eigenspace.
\end{proof}
\end{corollary}

\begin{lemma} \label{lem:adjoints-of-finitely-many-normal-matrices-polynomial} \hlt{Adjoints of Normal Matrices as One Polynomial}\\
Let $A_1,\ldots,A_k$ be normal complex matrices. Then $\exists$ a polynomial $p\in\mathbb{C}[z]$ s.t. $p(A_j)=A_j^* \ \forall j=1,\ldots,k$.
\begin{proof}
Let $\Omega=\bigcup\limits_{j=1}^{k}\operatorname{spec}(A_j)$. Since $\Omega$ is finite, choose a Lagrange interpolation polynomial $p$ satisfying $p(\lambda)=\overline{\lambda}$ for every $\lambda\in\Omega$. For each $A_j$, Corollary \ref{cor:normal-spectral-resolution-functional-calculus} gives $p(A_j)=A_j^*$.\\
\hlt{Structural Commentary}: One polynomial interpolates conjugation on all spectra.\\
\hlt{Key Insight}: For normal matrices, $A^*$ is $p(A)$ with $p(\lambda)=\overline{\lambda}$.
\end{proof}
\end{lemma}

\begin{theorem} \label{thm:fuglede-putnam-normal-intertwining} \hlt{Fuglede--Putnam Theorem}\\
Let $A\in\operatorname{Mat}(m,m,\mathbb{C})$ and $B\in\operatorname{Mat}(n,n,\mathbb{C})$ be normal, and let $X\in\operatorname{Mat}(m,n,\mathbb{C})$.\\
Then $AX=XB$ $\Longleftrightarrow$ $A^*X=XB^*$.
\begin{proof}
\begin{enumerate}[label=\roman*.]
\setlength{\itemsep}{0pt}
\item[$(\Rightarrow)$] Assume $AX=XB$. Then $A^rX=XB^r$ for every $r\geq0$, so $p(A)X=Xp(B)$ for every polynomial $p$. By Lemma \ref{lem:adjoints-of-finitely-many-normal-matrices-polynomial}, choose $p$ with $p(A)=A^*$ and $p(B)=B^*$. Thus $A^*X=XB^*$.
\item[$(\Leftarrow)$] Assume $A^*X=XB^*$. Apply the same argument to the normal matrices $A^*$ and $B^*$, choosing $q$ with $q(A^*)=A$ and $q(B^*)=B$. Hence $AX=XB$.
\end{enumerate}
\hlt{Structural Commentary}: An intertwining relation passes through every polynomial in the two normal matrices.\\
\hlt{Key Insight}: Normality makes adjoints polynomial, so intertwining transfers to adjoints.
\end{proof}
\end{theorem}

\begin{method} \label{meth:finding-spectral-decomposition-normal-matrix} \hlt{Finding a Spectral Decomposition of a Normal Matrix}\\
For a normal matrix $A\in\operatorname{Mat}(n,n,\mathbb{C})$:
\begin{enumerate}[label=\roman*.]
\setlength{\itemsep}{0pt}
\item Find the distinct eigenvalues $\mu_1,\ldots,\mu_d$.
\item For each $j$, solve $(A-\mu_jI_n)\mathbf{x}=0$ to get $E_{\mu_j}(A)$.
\item Orthonormalise a basis of each eigenspace.
\item Concatenate these orthonormal bases as the columns of $U$; then $U$ is unitary and $A=U\Lambda U^*$, where $\Lambda$ repeats each $\mu_j$ according to $\dim E_{\mu_j}(A)$.
\end{enumerate}
\end{method}

\begin{theorem} \label{thm:normal-spectral-classification-by-eigenvalues} \hlt{Spectral Classification of Normal Operators}\\
Let $\mathcal{V}$ be a finite-dimensional complex inner-product space and let $T\in\LL(\mathcal{V})$ be normal.
\begin{enumerate}[label=(\roman*)]
\setlength{\itemsep}{0pt}
\item $T$ is unitary if and only if $\abs{\lambda}=1$ for every $\lambda\in\operatorname{spec}(T)$.
\item $T$ is self-adjoint if and only if $\operatorname{spec}(T)\subseteq\mathbb{R}$.
\item $T$ is an orthogonal projection if and only if $\operatorname{spec}(T)\subseteq\{0,1\}$.
\item $T$ is skew-adjoint, meaning $T^*=-T$, if and only if every $\lambda\in\operatorname{spec}(T)$ is purely imaginary.
\end{enumerate}
\begin{proof}
Write the spectral resolution $T=\sum\limits_j\mu_jP_j$, where the $\mu_j$ are the distinct eigenvalues and $P_j\neq0$.\\
Then $I_{\mathcal{V}}=\sum_jP_j$, $T^*=\sum\limits_j\overline{\mu_j}P_j$, $T^*T=\sum\limits_j\abs{\mu_j}^2P_j$, and $T^2=\sum\limits_j\mu_j^2P_j$.\\
Since ranges of $P_j$ are nonzero and independent, an identity $\sum_j\alpha_jP_j=0$ holds only when every $\alpha_j=0$.
\begin{enumerate}[label=(\roman*)]
\setlength{\itemsep}{0pt}
\item $T$ is unitary iff $T^*T=I_{\mathcal{V}}=\sum_jP_j$, iff $\abs{\mu_j}=1$ for every $j$.
\item $T$ is self-adjoint iff $T=T^*$, iff $\mu_j=\overline{\mu_j}$ for every $j$.
\item $T$ is an orthogonal projection iff $T=T^*$ and $T^2=T$, iff $\mu_j=\overline{\mu_j}$ and $\mu_j^2=\mu_j \ \forall j$, iff $\mu_j\in\{0,1\} \ \forall j$.
\item $T^*=-T$ iff $\overline{\mu_j}=-\mu_j$ for every $j$, iff every $\mu_j$ is purely imaginary.
\end{enumerate}
\hlt{Structural Commentary}: For normal operators, structural classes are scalar conditions on spectral blocks.\\
\hlt{Key Insight}: For normal $T$, these properties reduce to scalar tests.
\end{proof}
\end{theorem}

\begin{proposition} \label{prop:cartesian-decomposition-normal-commuting-parts} \hlt{Cartesian Decomposition and Normality}\\
Let $\mathcal{V}$ be a finite-dimensional complex inner-product space.\\
For every $T\in\LL(\mathcal{V})$, there are unique self-adjoint operators $H,K\in\LL(\mathcal{V})$ such that $T=H+iK$.\\
They are $H=\frac{1}{2}(T+T^*)$ and $K=\frac{1}{2i}(T-T^*)$. Moreover, $T$ is normal if and only if $HK=KH$.
\begin{proof}
The displayed formulas give self-adjoint $H,K$ and $T=H+iK$. If $T=X+iY$ with $X,Y$ self-adjoint, then $T^*=X-iY$, so $X=\frac{1}{2}(T+T^*)$ and $Y=\frac{1}{2i}(T-T^*)$; uniqueness follows.\\
For $T=H+iK$, $T^*T=(H-iK)(H+iK)=H^2+K^2+i(HK-KH)$,\\
whereas $TT^*=(H+iK)(H-iK)=H^2+K^2-i(HK-KH)$. Thus $T^*T=TT^*$ if and only if $HK=KH$.\\
\hlt{Structural Commentary}: Every complex operator splits into Hermitian and skew-Hermitian parts.\\
\hlt{Key Insight}: Normality says the two Hermitian coordinates of $T$ commute.
\end{proof}
\end{proposition}

\begin{figure}[H]
  \centering
  \begin{tikzpicture}[
  >=Stealth,
  every node/.style={font=\small},
  box/.style={draw=gray!55, rounded corners=3pt, fill=gray!6, align=center, inner xsep=7pt, inner ysep=4.5pt, minimum height=0.9cm},
  result/.style={box, fill=green!7, minimum width=4.05cm},
  arrow/.style={->, line width=0.85pt, blue!70!black}
]
  \node[box, minimum width=1.4cm] (T) at (0,0) {$T$};
  \node[box, minimum width=3.75cm] (parts) at (2.95,0)
    {$H=\frac12(T+T^*)$\\[-1pt]
     $K=\frac{1}{2i}(T-T^*)$\\[-1pt]
     $T=H+iK$};
  \node[result] (normal) at (7.35,0)
    {$T\ \text{normal}\quad\Longleftrightarrow\quad HK=KH$};

  \draw[arrow] (T) -- (parts);
  \draw[arrow] (parts) -- (normal);
\end{tikzpicture}

  \caption{Cartesian decomposition: $T$ is normal exactly when $HK=KH$.}
  \label{fig:cartesian-normal-commutation}
\end{figure}

\begin{lemma} \label{lem:upper-triangular-normal-matrix-is-diagonal} \hlt{Upper-Triangular Normal Matrices Are Diagonal}\\
If $A\in\operatorname{Mat}(n,n,\mathbb{C})$ is upper triangular and normal, then $A$ is diagonal.
\begin{proof}
Induct on $n$. The case $n=1$ is immediate.\\
Since $A$ is upper triangular, $A\mathbf{e}_1=a_{11}\mathbf{e}_1$, while $A^*\mathbf{e}_1=\overline{a_{11}}\mathbf{e}_1+\overline{a_{12}}\mathbf{e}_2+\cdots+\overline{a_{1n}}\mathbf{e}_n$.\\
Normality gives $\norm{A\mathbf{e}_1}_2=\norm{A^*\mathbf{e}_1}_2$, hence $a_{12}=\cdots=a_{1n}=0$. Thus $A=a_{11}\oplus B$ with $B$ upper triangular. The block diagonal normality equation gives $B^*B=BB^*$, so $B$ is normal. By induction, $B$ is diagonal.\\
\hlt{Structural Commentary}: Normality forbids upper-triangular row tails.\\
\hlt{Key Insight}: Compare the lengths of $A\mathbf{e}_1$ and $A^*\mathbf{e}_1$ to kill the first row, then induct.
\end{proof}
\end{lemma}

\begin{proposition} \label{prop:block-upper-triangular-normal-collapse} \hlt{Normal Block Upper-Triangular Matrices Collapse}\\
Let $\mathbb{F}\in\{\mathbb{R},\mathbb{C}\}$. Let $A=
\begin{bmatrix}
B&X\\
0&C
\end{bmatrix}
\in\operatorname{Mat}(n,n,\mathbb{F})$, where $B$ and $C$ are square.\\
Then $A$ is normal if and only if $X=0$ and both $B$ and $C$ are normal.
\begin{proof}
\begin{enumerate}[label=\roman*.]
\setlength{\itemsep}{0pt}
\item[$(\Rightarrow)$] Suppose $A$ is normal. Comparing the $(1,1)$ blocks of $A^*A$ and $AA^*$ gives $B^*B=BB^*+XX^*$.\\
Taking traces and using cyclicity, $\operatorname{tr}(XX^*)=\operatorname{tr}(B^*B)-\operatorname{tr}(BB^*)=0$. Since $\operatorname{tr}(XX^*)=\norm{X}_F^2$, we get $X=0$. Thus $A=B\oplus C$, and $A^*A=AA^*$ reduces blockwise to $B^*B=BB^*$ and $C^*C=CC^*$.
\item[$(\Leftarrow)$] If $X=0$ and $B,C$ are normal, then $A=B\oplus C$, so blockwise multiplication gives $A^*A=AA^*$.
\end{enumerate}
\hlt{Structural Commentary}: A normal operator cannot have one-way coupling between orthogonal blocks.\\
\hlt{Key Insight}: The off-diagonal block contributes the nonnegative trace $\norm{X}_F^2$, so normality forces it to vanish.
\end{proof}
\end{proposition}

\begin{figure}[H]
  \centering
  \begin{tikzpicture}[
  >=Stealth,
  every node/.style={font=\small},
  block/.style={draw=gray!55, rounded corners=3pt, inner xsep=6pt, inner ysep=5pt, align=center},
  arrow/.style={->, line width=0.85pt, blue!70!black}
]
  \node[block] (tri) at (0,0)
    {$A=\begin{bmatrix}B&X\\0&C\end{bmatrix}$\\normal};
  \node[block] (schur) at (3.55,0)
    {$T_B=U_B^*BU_B,\quad T_C=U_C^*CU_C$\\
     $Q^*AQ=\begin{bmatrix}T_B&U_B^*XU_C\\0&T_C\end{bmatrix}$\\
     upper triangular and normal};
  \node[block] (diag) at (7.62,0)
    {$Q^*AQ$ diagonal\\$X=0$, $B,C$ normal};

  \draw[arrow] (tri) -- (schur);
  \draw[arrow] (schur) -- (diag);
\end{tikzpicture}

  \caption{Normality removes the upper-right block in a block upper-triangular matrix.}
  \label{fig:block-upper-triangular-normal-collapse}
\end{figure}

\begin{theorem} \label{thm:normal-invariant-subspaces-reduce} \hlt{Invariant Subspaces of Normal Matrices Reduce}\\
Let $\mathbb{F}\in\{\mathbb{R},\mathbb{C}\}$, let $A\in\operatorname{Mat}(n,n,\mathbb{F})$ be normal, and let $\mathcal{U}\subseteq\mathbb{F}^n$ be invariant under $A$. Then $\mathcal{U}$ is invariant under $A^*$. If $P$ is the orthogonal projection onto $\mathcal{U}$, then $AP=PA$.
\begin{proof}
The cases $\mathcal{U}=\{0\}$ and $\mathcal{U}=\mathbb{F}^n$ are immediate. Otherwise choose a unitary $Q=[Q_1\ Q_2]$ whose columns form orthonormal bases of $\mathcal{U}$ and $\mathcal{U}^{\perp}$. Invariance of $\mathcal{U}$ says $AQ_1$ has columns in $\mathcal{U}=\operatorname{col}(Q_1)$, hence $Q_2^*AQ_1=0$.\\
Therefore $Q^*AQ=
\begin{bmatrix}
B&X\\
0&C
\end{bmatrix}$. This matrix is normal because it is unitarily similar to $A$.\\
Proposition \ref{prop:block-upper-triangular-normal-collapse} gives $X=0$, so $Q^*AQ=B\oplus C$. Hence $Q^*A^*Q=B^*\oplus C^*$, which shows $\mathcal{U}$ is $A^*$-invariant.\\
Also $Q^*PQ=I_{\dim(\mathcal{U})}\oplus0$, so $Q^*AQ$ commutes with $Q^*PQ$; conjugating back gives $AP=PA$.\\
\hlt{Structural Commentary}: In adapted coordinates, normal invariance becomes block diagonality.\\
\hlt{Key Insight}: Normal invariant subspaces are automatically reducing subspaces.
\end{proof}
\end{theorem}

\begin{corollary} \label{cor:normal-operator-invariant-subspaces-reduce} \hlt{Invariant Subspaces of Normal Operators Reduce}\\
Let $\mathcal{V}$ be a finite-dimensional inner-product space, let $T\in\LL(\mathcal{V})$ be normal, let $\mathcal{U}$ be invariant under $T$. Then:
\begin{enumerate}[label=(\roman*)]
\setlength{\itemsep}{0pt}
\item $\mathcal{U}^{\perp}$ is invariant under $T$;
\item $\mathcal{U}$ is invariant under $T^*$;
\item $(T|_{\mathcal{U}})^*=(T^*)|_{\mathcal{U}}$;
\item $T|_{\mathcal{U}}$ and $T|_{\mathcal{U}^{\perp}}$ are normal.
\end{enumerate}
\begin{proof}
Choose an orthonormal basis adapted to $\mathcal{U}\oplus\mathcal{U}^{\perp}$.\\
In these coordinates, Theorem \ref{thm:normal-invariant-subspaces-reduce} makes the matrix of $T$ block diagonal, say $A=B\oplus C$.\\
The block diagonal form shows $\mathcal{U}^{\perp}$ is invariant under $T$ and $\mathcal{U}$ is invariant under $T^*$. It also gives $(T|_{\mathcal{U}})^*=(T^*)|_{\mathcal{U}}$ because the adjoint block on $\mathcal{U}$ is $B^*$. Finally, normality of $A$ reduces blockwise to normality of $B$ and $C$.\\
\hlt{Structural Commentary}: Normal operator turns every invariant subspace into an orthogonal reducing summand.\\
\hlt{Key Insight}: Normality kills the one-way coupling between $\mathcal{U}$ and $\mathcal{U}^{\perp}$.
\end{proof}
\end{corollary}

\begin{theorem} \label{thm:simultaneous-unitary-diagonalization-normal-family} \hlt{Simultaneous Unitary Diagonalisation of Commuting Normal Operators}\\
Let $\emptyset\neq\mathcal{F}\subseteq\LL(\mathcal{V})$ be a normal family on a finite-dimensional complex inner-product space. Then $\mathcal{F}$ is pairwise commuting if and only if $\mathcal{V}$ has an orthonormal basis of common eigenvectors for $\mathcal{F}$.
\begin{proof}
\begin{enumerate}[label=\roman*.]
\setlength{\itemsep}{0pt}
\item[$(\Leftarrow)$] In a common orthonormal eigenbasis, every $A\in\mathcal{F}$ is diagonal. Hence all operators in $\mathcal{F}$ commute.
\item[$(\Rightarrow)$] Induct on $\dim(\mathcal{V})$. The zero-dimensional case is vacuous. If every $A\in\mathcal{F}$ is scalar, any orthonormal basis works. Otherwise choose a nonscalar $A_0\in\mathcal{F}$ and write its spectral resolution as $A_0=\sum_j\mu_jP_j$.\\
Since $A_0$ is nonscalar, at least two nonzero ranges $\operatorname{range}(P_j)$ occur, and each has dimension strictly smaller than $\dim(\mathcal{V})$. By Corollary \ref{cor:normal-spectral-resolution-functional-calculus}, each $P_j$ is a polynomial in $A_0$.\\
If $B\in\mathcal{F}$, then $BA_0=A_0B$, so $B$ commutes with every polynomial in $A_0$, in particular with $P_j$. Hence $B(\operatorname{range}(P_j))\subseteq\operatorname{range}(P_j)$ because $P_jB\mathbf{x}=BP_j\mathbf{x}=B\mathbf{x}$ whenever $P_j\mathbf{x}=\mathbf{x}$.\\
The same commutation preserves $\ker(P_j)=\operatorname{range}(P_j)^{\perp}$.\\
In the decomposition $\mathcal{V}=\operatorname{range}(P_j)\oplus\ker(P_j)$, the matrix of $B$ is therefore block diagonal; since $B$ is normal, its diagonal restrictions are normal block by block.\\
Thus, on each nonzero $\operatorname{range}(P_j)$, the restricted family is again pairwise commuting and normal.\\
By induction it has a common orthonormal eigenbasis. Concatenating these bases over all $j$ gives a common orthonormal eigenbasis of $\mathcal{V}$.
\end{enumerate}
\hlt{Structural Commentary}: Spectral projections reduce the whole commuting family.\\
\hlt{Key Insight}: Diagonalise one nonscalar member, then induct on its eigenspaces.
\end{proof}
\end{theorem}

\begin{figure}[H]
  \centering
  \begin{tikzpicture}[
  >=Stealth,
  every node/.style={font=\small},
  block/.style={draw=gray!55, rounded corners=3pt, inner xsep=8pt, inner ysep=4.5pt, align=center},
  arrow/.style={->, line width=0.8pt, blue!70!black, shorten >=2pt, shorten <=2pt}
]
  \node[block] (family) at (0,0) {$\mathcal{F}$};
  \node[block] (split) at (2.46,0) {$A_0=\sum_j\mu_jP_j$};
  \node[block] (pieces) at (5.32,0) {$\bigoplus_j\operatorname{range}(P_j)$};
  \node[block] (basis) at (8.49,0) {common eigenbasis};

  \draw[arrow] (family.east) -- (split.west);
  \draw[arrow] (split.east) -- (pieces.west);
  \draw[arrow] (pieces.east) -- (basis.west);
\end{tikzpicture}

  \caption{For $A_0=\sum_j\mu_jP_j$, commutativity gives $BP_j=P_jB$; the ranges reduce $\mathcal{F}$, and induction gives a common orthonormal eigenbasis.}
  \label{fig:simultaneous-normal-diagonalization}
\end{figure}

\begin{corollary} \label{cor:commuting-normal-sums-products-normal} \hlt{Sums and Products of Commuting Normal Operators}\\
Let $S,T\in\LL(\mathcal{V})$ be commuting normal operators on a finite-dimensional complex inner-product space.\\
Then $S+T$ and $ST$ are normal.
\begin{proof}
By Theorem \ref{thm:simultaneous-unitary-diagonalization-normal-family}, $S$ and $T$ are diagonal in a common orthonormal basis.\\
Then $S+T$ and $ST$ are also diagonal in that basis, hence normal.\\
\hlt{Structural Commentary}: Commuting normal operators share orthonormal coordinates.\\
\hlt{Key Insight}: In common diagonal form, sums and products remain diagonal.
\end{proof}
\end{corollary}

\begin{theorem} \label{thm:normality-transfer-ab-ba} \hlt{Normality Transfers from $AB$ to $BA$}\\
Let $A,B\in\operatorname{Mat}(n,n,\C)$. If $A$, $B$, and $AB$ are normal, then $BA$ is normal.
\begin{proof}
By Theorem \ref{thm:ab-ba-nonzero-eigenvalues}, $AB$ and $BA$ have the same eigenvalues counted with algebraic multiplicity. Also,
\begin{align}
\operatorname{tr}\bigl((BA)^*(BA)\bigr) &=\operatorname{tr}(A^*B^*BA) =\operatorname{tr}(AA^*B^*B) \nonumber\\
&=\operatorname{tr}(A^*AB^*B) =\operatorname{tr}(A^*A BB^*) =\operatorname{tr}(B^*A^*AB) =\operatorname{tr}\bigl((AB)^*(AB)\bigr),
\nonumber
\end{align}
using cyclicity of trace and normality of $A$ and $B$. Thus $\norm{BA}_F^2=\norm{AB}_F^2$, so $\Delta(BA)=\Delta(AB)$.\\
Since $AB$ is normal, $\Delta(AB)=0$; hence $\Delta(BA)=0$, and Theorem \ref{thm:schur-inequality-normal-defect} gives that $BA$ is normal.\\
\hlt{Structural Commentary}: Same spectrum and Frobenius norm give the same Schur defect.\\
\hlt{Key Insight}: Schur defect turns normality of $AB$ into normality of $BA$.
\end{proof}
\end{theorem}

\begin{definition} \label{def:circulant-matrix-cyclic-shift} \hlt{Circulant Matrix}\\
On the $0$-indexed basis of $\C^n$, define the \hlt{cyclic shift}
$P_n\mathbf{e}_j=\mathbf{e}_{(j-1)\operatorname{mod} n}$.\\
A matrix $C\in\operatorname{Mat}(n,n,\C)$ is \hlt{circulant} if, for some $a_0,\ldots,a_{n-1}\in\C$,
\begin{align}
C&=a_0I_n+a_1P_n+\cdots+a_{n-1}P_n^{n-1}=p(P_n), \nonumber\\
p(z)&=a_0+a_1z+\cdots+a_{n-1}z^{n-1} \nonumber
\end{align}
\end{definition}

\begin{theorem} \label{thm:fourier-basis-diagonalizes-cyclic-circulants} \hlt{Fourier Basis for Cyclic Shifts and Circulants}\\
Let $n\geq1$, let $\omega_n\defeq e^{2\pi i/n}$, and let $P_n$ be the cyclic shift.\\
For $0\leq k<n$, set $\mathbf{f}_k\defeq(\omega_n^{jk})_{j=0}^{n-1}$; $F_n\defeq[\mathbf{f}_0\ \cdots\ \mathbf{f}_{n-1}]$; and $\Omega_n\defeq\operatorname{diag}(1,\omega_n,\ldots,\omega_n^{n-1})$.\\
Then $F_n^*F_n=nI_n$, so $Q_n\defeq n^{-1/2}F_n$ is unitary, and $P_nF_n=F_n\Omega_n$.\\
If $C=p(P_n)$ is circulant, then $Q_n^*CQ_n=p(\Omega_n) =\operatorname{diag}\bigl(p(1),p(\omega_n),\ldots,p(\omega_n^{n-1})\bigr)$.
\begin{proof}
For $0\leq k,\ell\leq n-1$, $\innerprod{\mathbf{f}_k,\mathbf{f}_\ell} =\sum\limits_{r=0}^{n-1}(\omega_n^{k-\ell})^r =n\delta_{k\ell}$.
For $k=\ell$ the sum is $n$; for $k\neq\ell$, $W=\omega_n^{k-\ell}$ has $W\neq1$ and $W^n=1$, so $(1-W)\sum\limits_{r=0}^{n-1}W^r=0$. Hence $F_n^*F_n=nI_n$, and Theorem \ref{thm:unitary-matrix-equivalent-conditions} makes $Q_n$ unitary.\\
Also $P_n\mathbf{f}_k=\omega_n^k\mathbf{f}_k$, so $Q_n^*P_nQ_n=\Omega_n$; applying $p$ gives the formula.\\
\hlt{Structural Commentary}: Fourier coordinates diagonalize every circulant.\\
\hlt{Key Insight}: Orthogonality of roots of unity turns a cyclic matrix problem into scalar evaluation.
\end{proof}
\end{theorem}

\begin{figure}[H]
  \centering
  \begin{tikzpicture}[
  >=Stealth,
  x=0.92cm,
  y=0.9cm,
  every node/.style={font=\small},
  box/.style={draw=gray!55, rounded corners=3pt, fill=gray!6, align=center, minimum width=2.15cm, minimum height=0.8cm},
  result/.style={box, fill=green!7},
  arrow/.style={->, line width=0.85pt, blue!70!black},
  lbl/.style={fill=white, text=gray!70!black, inner xsep=2pt, inner ysep=1pt}
]
  \node[box] (shift) at (0,1.1) {$P_n$};
  \node[box] (diagshift) at (4.15,1.1) {$\Omega_n$};
  \node[box] (circ) at (0,-1.1) {$C=p(P_n)$};
  \node[result] (diagcirc) at (4.15,-1.1) {$p(\Omega_n)$};

  \draw[arrow] (shift) -- node[lbl, above] {$Q_n^*(\,\cdot\,)Q_n$} (diagshift);
  \draw[arrow] (circ) -- node[lbl, below] {$Q_n^*(\,\cdot\,)Q_n$} (diagcirc);
  \draw[arrow] (shift) -- node[lbl, left] {$p(\cdot)$} (circ);
  \draw[arrow] (diagshift) -- node[lbl, right] {$p(\cdot)$} (diagcirc);
\end{tikzpicture}

  \caption{Fourier diagonalizes $P_n$; polynomial evaluation diagonalizes $C=p(P_n)$.}
  \label{fig:circulant-fourier-diagonalization}
\end{figure}

\begin{proposition} \label{prop:permutation-similarity-reorders-diagonal} \hlt{Permutation Similarity Reorders Diagonal Entries}\\
Let $\Lambda=\operatorname{diag}(\lambda_1,\ldots,\lambda_n)$ and let $P=[\mathbf{e}_{\sigma(1)}\ \cdots\ \mathbf{e}_{\sigma(n)}]$ be a permutation matrix. Then $P^*\Lambda P=\operatorname{diag}(\lambda_{\sigma(1)},\ldots,\lambda_{\sigma(n)})$.
\begin{proof}
The $(i,j)$ entry of $P^*\Lambda P$ is $\mathbf{e}_{\sigma(i)}^*\Lambda\mathbf{e}_{\sigma(j)} =\lambda_{\sigma(j)}\mathbf{e}_{\sigma(i)}^*\mathbf{e}_{\sigma(j)} =\lambda_{\sigma(j)}\delta_{ij}$.\\
Thus all off-diagonal entries are zero, and the $j$th diagonal entry is $\lambda_{\sigma(j)}$.\\
\hlt{Structural Commentary}: Permutation similarity applies the same basis reordering to rows and columns.\\
\hlt{Key Insight}: Diagonal entries move exactly as the standard basis vectors are rearranged.
\end{proof}
\end{proposition}

\begin{corollary} \label{cor:normal-matrices-same-spectrum-unitarily-similar} \hlt{Normal Matrices with the Same Spectrum Are Unitarily Similar}\\
If $A,B\in\operatorname{Mat}(n,n,\mathbb{C})$ are normal and have the same eigenvalues counted with algebraic multiplicity, then $A$ and $B$ are unitarily similar.
\begin{proof}
By Theorem \ref{thm:complex-spectral-theorem}, write $A=U\Lambda U^*$ and $B=V\Gamma V^*$ with $U,V$ unitary and $\Lambda,\Gamma$ diagonal.\\
Since eigenvalues agree with multiplicity, Proposition \ref{prop:permutation-similarity-reorders-diagonal} gives a permutation matrix $P$ such that $\Gamma=P^*\Lambda P$.\\
Hence $B=(VP^*U^*)A(UPV^*)$, and $UPV^*$ is unitary.\\
\hlt{Structural Commentary}: Normal matrices are classified up to unitary similarity by their eigenvalue multiset.\\
\hlt{Key Insight}: After diagonalising both matrices, only a permutation of diagonal entries remains.
\end{proof}
\end{corollary}

\begin{theorem} \label{thm:complex-isometry-description} \hlt{Description of Complex Isometries}\\
Let $\mathcal{V}$ be a finite-dimensional complex inner-product space and let $S\in\LL(\mathcal{V})$. Then $S$ is an isometry if and only if $\mathcal{V}$ has an orthonormal basis consisting of eigenvectors of $S$ whose eigenvalues all have absolute value $1$.
\begin{proof}
If $S\mathbf{e}_j=\lambda_j\mathbf{e}_j$ in an orthonormal basis with $\abs{\lambda_j}=1$, then for $\mathbf{v}=\sum_j c_j\mathbf{e}_j$,\\
$\norm{S\mathbf{v}}^2=\sum\limits_j\abs{\lambda_jc_j}^2=\sum\limits_j\abs{c_j}^2=\norm{\mathbf{v}}^2$, so $S$ is an isometry. Conversely, if $S$ is an isometry, then $S^*S=I_{\mathcal{V}}$ by Theorem \ref{thm:finite-dimensional-isometry-iff-adjoint-product-identity}. Finite-dimensional invertibility gives $SS^*=I_{\mathcal{V}}$, so $S$ is normal. By Complex Spectral Theorem, choose orthonormal eigenbasis $S\mathbf{e}_j=\lambda_j\mathbf{e}_j$. Since $\mathbf{e}_j$ is unit and $S$ preserves norms, $\abs{\lambda_j}=\norm{S\mathbf{e}_j}=1$.\\
\hlt{Structural Commentary}: Complex isometries are normal with unit-circle spectrum.\\
\hlt{Key Insight}: In an eigenbasis, isometry means $\abs{\lambda_j}=1$.
\end{proof}
\end{theorem}

\begin{definition} \label{def:positive-operator} \hlt{Positive Operator}\\
$T\in\LL(\mathcal{V})$ is \hlt{positive} if $T=T^*$ and $\innerprod{T\mathbf{v},\mathbf{v}}\geq0$ for all $\mathbf{v}\in\mathcal{V}$.
\end{definition}

\begin{definition} \label{def:square-root-operator} \hlt{Square Root of an Operator}\\
An operator $R\in\LL(\mathcal{V})$ is a \hlt{square root} of $T\in\LL(\mathcal{V})$ if $R^2=T$.
\end{definition}

\begin{theorem} \label{thm:positive-operator-characterization} \hlt{Characterization of Positive Operators}\\
For $T\in\LL(\mathcal{V})$, the following are equivalent:
\begin{enumerate}[label=\roman*.]
\setlength{\itemsep}{0pt}
\item $T$ is positive;
\item $T$ is self-adjoint and every eigenvalue of $T$ is nonnegative;
\item $T$ has a positive square root;
\item $T$ has a self-adjoint square root;
\item there exists $R\in\LL(\mathcal{V})$ such that $T=R^*R$.
\end{enumerate}
\begin{proof}
\begin{enumerate}[label=(\roman*).]
\setlength{\itemsep}{0pt}
\item[$(i)$]$\Rightarrow (ii)$ If $T\mathbf{v}=\lambda\mathbf{v}$ and $\mathbf{v}\neq\mathbf{0}$, then $\lambda\in\mathbb{R}$ by Theorem \ref{thm:self-adjoint-eigenvalues-real}, and $0\leq\innerprod{T\mathbf{v},\mathbf{v}}=\lambda\norm{\mathbf{v}}^2$, so $\lambda\geq0$.
\item[$(ii)$]$\Rightarrow (iii)$ By the real or complex spectral theorem, choose an orthonormal eigenbasis $(\mathbf{e}_1,\ldots,\mathbf{e}_n)$ with $T\mathbf{e}_j=\lambda_j\mathbf{e}_j$ and $\lambda_j\geq0$. Define $S\mathbf{e}_j=\sqrt{\lambda_j}\mathbf{e}_j$ and extend linearly.\\
Then $S$ is diagonal with nonnegative real diagonal entries in an orthonormal basis, hence self-adjoint and positive, and $S^2\mathbf{e}_j=\lambda_j\mathbf{e}_j=T\mathbf{e}_j$ for every $j$. Thus $S^2=T$.
\item[$(iii)$]$\Rightarrow (iv)$ A positive square root is self-adjoint.
\item[$(iv)$]$\Rightarrow (v)$ If $T=R^2$ with $R=R^*$, then $T=R^*R$.
\item[$(v)$]$\Rightarrow (i)$ If $T=R^*R$, then $T^*=T$ and $\innerprod{T\mathbf{v},\mathbf{v}}=\innerprod{R\mathbf{v},R\mathbf{v}}=\norm{R\mathbf{v}}^2\geq0$ for every $\mathbf{v}$, so $T$ is positive.
\end{enumerate}
\hlt{Structural Commentary}: Positivity equals nonnegative spectrum and square-root factorization.\\
\hlt{Key Insight}: In an eigenbasis, the proof is scalar: $\lambda\geq0$ iff $\sqrt{\lambda}$ exists.
\end{proof}
\end{theorem}

\begin{proposition} \label{prop:positive-square-root-unique} \hlt{Unique Positive Square Root}\\
Every positive operator on $\mathcal{V}$ has a unique positive square root.
\begin{proof}
Existence follows from Theorem \ref{thm:positive-operator-characterization}. For uniqueness, decompose $\mathcal{V}=E_{\lambda_1}(T)\oplus\cdots\oplus E_{\lambda_m}(T)$.\\
Let $R$ be a positive square root of $T$. Since $R^2=T$, we have $RT=R^3=TR$, so each $E_{\lambda_j}(T)$ is invariant under $R$. On $E_{\lambda_j}(T)$, the restriction $R_j\defeq R|_{E_{\lambda_j}(T)}$ is positive and satisfies $R_j^2=\lambda_jI$.\\
By the spectral theorem for the positive operator $R_j$, every eigenvalue $\rho$ of $R_j$ satisfies $\rho\geq0$ and $\rho^2=\lambda_j$, hence $\rho=\sqrt{\lambda_j}$. Therefore $R_j=\sqrt{\lambda_j}I$ on $E_{\lambda_j}(T)$.\\
Thus every positive square root agrees with the same scalar operator on each eigenspace $E_{\lambda_j}(T)$, so it is unique.\\
\hlt{Structural Commentary}: Each eigenspace determines its own scalar root.\\
\hlt{Key Insight}: Positivity selects the nonnegative root.
\end{proof}
\end{proposition}

\begin{definition} \label{def:positive-square-root-notation} \hlt{Positive Square Root Notation}\\
If $T$ is positive, then $\sqrt{T}$ denotes the unique positive square root of $T$.
\end{definition}
